\begin{table}[htbp!]
  \caption{Targets}
  \label{tab:targets}
  \centering
  \begin{small}
  \resizebox{0.8\textwidth}{!}{\begin{tabular}{llll}
    \toprule
    Target & Symbol & \multicolumn{2}{c}{Value} \\
           &        & Home & Foreign \\
    \midrule
    Relative Reaction of Tradable to Non-Tradable Employment &  & 0.027 & -- \\
    Non-Tradable Employment Response &  & -0.0060 & -- \\
    Unemployment Rate & $u$ & 0.089 & 0.096 \\
    Job-Finding Rate & $\rho$ & 0.116 & 0.108 \\
    Tradable Producer Price & $p^T_H$ & 1 & 1 \\
    CPI & $P$ & 1 & 1 \\
    Real Exchange Rate & $\mathrm{RER}$ & 1 & 1 \\
    Debt-to-GDP Ratio & $b/Y$ & 0.600 & 0.600 \\
    Pre-Reform Net Foreign Assets & $NFA/Y$ & 0 & 0 \\
    \bottomrule
  \end{tabular}}
  \end{small}
  \begin{flushleft}
  \footnotesize \textit{Notes:} Steady-state and transition targets of the calibration. The sectoral employment responses are averages over the first 40 quarters and refer to the full reform, that is, both the replacement-rate cut and the entitlement-duration cut. They are matched jointly by the replacement-rate cut and the non-tradable recruiting curvature $\chi_N$. The zero pre-reform net foreign asset position is matched by the foreign discount factor. The job-filling rate $q=0.7$ is used as a normalization of the vacancy scale and is reported in Table~\ref{tab:calib}.
  \end{flushleft}
\end{table}
